%===============================================================================================================%
%                    Estilo y paquetes que se van a utilizar: No es necesario editar                            %
%===============================================================================================================%
\documentclass[12pt,letterpaper]{article}     % Tipo de documento y otras especificaciones
\usepackage[utf8]{inputenc}                   % Para escribir tildes y eñes
\usepackage[spanish]{babel}                   % Para que los títulos de figuras, tablas y otros estén en español
\addto\captionsspanish{\renewcommand{\tablename}{Tabla}}					% Cambiar nombre a tablas
\addto\captionsspanish{\renewcommand{\listtablename}{Índice de tablas}}		% Cambiar nombre a lista de tablas
\usepackage{geometry}                          
\geometry{left=18mm,right=18mm,top=21mm,bottom=21mm}	 % Tamaño del área de escritura de la página
\usepackage{ucs}
\usepackage{amsmath}      % Los paquetes ams son desarrollados por la American Mathematical Society
\usepackage{amsfonts}     % y mejoran la escritura de fórmulas y símbolos matemáticos.
\usepackage{tikz}

\usetikzlibrary{babel}               % Se añade para que se acepten algunos caractéres en el circuitikz 
\usepackage{circuitikz}
\usepackage{graphicx}     % Para insertar gráficas
\usepackage[lofdepth,lotdepth]{subfig}	% Para colocar varias figuras
\usepackage{unitsdef}	  % Para la presentación correcta de unidades
\usepackage{pdfpages}     % Incluir paginas de pdf externo, para los anexos
\usepackage{appendix}     % Para los anexos
\renewcommand{\unitvaluesep}{\hspace*{4pt}}	% Redimensionamiento del espacio entre magnitud y unidad
\usepackage[colorlinks=true,urlcolor=blue,linkcolor=black,citecolor=red]{hyperref}     % Para insertar marcadores
\usepackage{float}		% Para ubicar las tablas y figuras justo después del texto
\usepackage{booktabs}	% Para hacer tablas más estilizadas
\batchmode
\usepackage{apalike}
\bibliographystyle{apalike}  % Bibliografía en formato IEEE
\pagestyle{plain} 
\pagenumbering{arabic}
\usepackage{lastpage}
\usepackage{skull}
\usepackage{marvosym}
\usepackage{fancyhdr}	% Para manejar los encabezados y pies de página
\pagestyle{fancy}		% Contenido de los encabezados y pies de pagina
\usepackage{multicol}   % Para varias columnas
\usepackage{enumerate}
\usepackage{url}
\usepackage{hyperref}
\usepackage{listings}
\usepackage{xcolor}

\lstset{
  language=C++,
  basicstyle=\ttfamily\footnotesize,
  keywordstyle=\color{blue}\bfseries,
  commentstyle=\color{gray}\ttfamily,
  stringstyle=\color{red!70!black},
  numbers=left,
  numberstyle=\tiny\color{gray},
  stepnumber=1,
  numbersep=5pt,
  backgroundcolor=\color{gray!10},
  frame=single,
  rulecolor=\color{black},
  captionpos=b,
  breaklines=true,
  breakatwhitespace=false,
  tabsize=2,
  showstringspaces=false
}

%===============================================================================================================%
%                         Definición del ambiente: Resumen                                                      %   
%===============================================================================================================%
\newcounter{resumen}
\setcounter{resumen}{0}
\def\theejemplo{\thechapter.\arabic{resumen}}

\newenvironment{resumen}
{	
	\begin{center}
	\begin{minipage}[t]{500 pt}
	\vspace{1mm}
	\emph{\textbf{}}
	\\[-2mm]
	\line(1,0){500}
	\\[-4.25 mm]
	\line(1,0){500}
	\\
}
{
	\normalsize
	\\[2mm]
	\footnotesize\textbf{Palabras clave: \footnotesize\@palabras}
	\\[-2mm]
	\line(1,0){500}
	\\[0.5cm]
	\end{minipage}
	\end{center}
}

%===============================================================================================================%
%                         Definición del ambiente: Palabras clave                                               %
%===============================================================================================================%
\def\palabras#1{\gdef\@palabras{#1}}

%===============================================================================================================%
%             Definición de encabezados                                                                         %
%===============================================================================================================%
\lhead{IoT - Arquitectura de Sistemas} 
\chead{}
\rhead{Manual Técnico}	
\lfoot{Programa de Ingeniería Electrónica y Sistemas}
\cfoot{\thepage\ de \pageref{LastPage}}
\rfoot{Universidad Católica de Oriente}

%===============================================================================================================%
% Portada                                                                                                       %
%===============================================================================================================%
\author{Santiago Barbarán, 1.001.159.841\\Miguel Angel Botero, 1.023.523.241\\Sebastian Moreno, 1.001.724.352 \\Profesor: Juan Camilo Acevedo \vspace*{2.8in}}
\title{
\begin{minipage}{\textwidth}
\centering
	\includegraphics[width=0.3\textwidth]{uco.png}
	\bigskip
\end{minipage}\\
Universidad Católica de Oriente \\{\small Facultad de Ingeniería\\Programa de Ingeniería Electrónica y Sistemas\\ Manual Técnico\\Semestre 2026-II\\\vspace*{0.55in} Documentación Técnica:}\\ Manual Tecnico\vspace*{1.35in}}

\begin{document}	% Inicio del documento

\pdfbookmark[1]{Portada}{portada} 	% Marcador para el título

\maketitle							% Título
\thispagestyle{empty}

\tableofcontents         % Tabla de contenidos
\newpage

\palabras{IoT, RP2040, Raspberry Pi 5, MQTT, TLS X.509, Mosquitto, Grafana, Electroválvula, Contactor.}

\begin{resumen}
Este manual técnico documenta el diseño, la arquitectura hardware/software y el protocolo de despliegue del sistema IoT para monitoreo y control automatizado de piscicultura. La arquitectura está dividida en cinco capas bien definidas: Capa de Percepción (sensores de conductividad/SDT y nivel JSN-SR04T), Capa de Nodos Perimetrales (microcontrolador RP2040 con lógica de control local), Capa de Red y Procesamiento (Gateway local Raspberry Pi 5 ejecutando broker MQTT Mosquitto), Capa Nube (transmisión segura bajo TLS/X.509) y Capa de Aplicación (Dashboard Grafana/Web Application para visualización en tiempo real). Se detallan los esquemáticos electrónicos, aislamiento galvánico de cargas de 110V AC mediante optoacopladores y relés/contactores, además de las configuraciones de software y scripts de ingesta.
\end{resumen}

\section{Introducción}
El monitoreo continuo de parámetros fisicoquímicos e hidráulicos en tanques de piscicultura es fundamental para mitigar la mortalidad de peces y optimizar la calidad del agua. El presente sistema IoT integra un esquema distribuido donde nodos de campo basados en la plataforma Raspberry Pi Pico (RP2040) adquieren variables de conductividad y nivel, ejecutando el control local de actuadores de 110V AC (bombas de recirculación y electroválvulas) mientras transmiten datos hacia un Gateway local centralizado basado en Raspberry Pi 5.

\section{Arquitectura del Sistema}
La arquitectura tecnológica del sistema sigue un modelo multicapa estructurado de la siguiente forma:

\subsection{Capa de Percepción / Campo}
\begin{itemize}
    \item \textbf{Sensor Ultrasónico de Nivel JSN-SR04T:} Módulo impermeable que mide la distancia al nivel libre de agua del estanque mediante interfaz serial UART/GPIO.
    \item \textbf{Conductivímetro / Sensor de SDT:} Módulo analógico equipado con electrodos que entrega una señal de voltaje proporcional a la concentración de Sólidos Disueltos Totales en el agua, conectada a los canales ADC del RP2040.
\end{itemize}

\subsection{Capa de Nodos Perimetrales}
\begin{itemize}
    \item \textbf{Microcontrolador RP2040 (Raspberry Pi Pico):} Procesa localmente las lecturas de los sensores, aplica filtros digitales y ejecuta decisiones de control local autónomo (activación de electroválvulas de drenaje/ingreso y bombas de recirculación).
\end{itemize}

\subsection{Capa de Red y Procesamiento (Gateway Local)}
\begin{itemize}
    \item \textbf{Raspberry Pi 5:} Actúa como servidor y gateway local, hospedando el Broker MQTT Mosquitto con mecanismos de autenticación y cifrado TLS/X.509.
    \item \textbf{Servicios locales:} Ejecución de base de datos de series temporales y motor de ingesta de datos.
\end{itemize}

\subsection{Capa Nube y Aplicación}
\begin{itemize}
    \item \textbf{Dashboard Grafana / Web Application:} Interfaz de usuario para visualización gráfica de métricas en tiempo real, histórico de parámetros y monitoreo del estado de actuadores.
\end{itemize}

\section{Esquemático y Etapa de Potencia}
Para evitar el acoplamiento de ruido inductivo procedente de la bomba de recirculación de 110V AC y las electroválvulas hacia el microcontrolador RP2040, se implementa una etapa de aislamiento optoelectrónico basada en optoacopladores MOC3021 o relés optoacoplados de estado sólido.

\begin{figure}[H]
\centering
\begin{circuitikz}[scale=1, transform shape]
    % Microcontrolador RP2040
    \draw[thick, fill=blue!5] (0,0) rectangle (5.5,5.2);
    \node at (2.8,4.7) {\textbf{RP2040 (Nodo Sensor)}};
    \node[right] at (0.2,4) {\footnotesize GP15 $\rightarrow$ TRIG (JSN-SR04T)};
    \node[right] at (0.2,3.0) {\footnotesize GP14 $\leftarrow$ ECHO (JSN-SR04T)};
    \node[right] at (0.2,2.2) {\footnotesize ADC0/GP26 $\leftarrow$ Sensor SDT};
    \node[left] at (4.3,1.4) {\footnotesize GP16 (Carga TDS)};
    \node[left] at (4.3,0.6) {\footnotesize GP17 (Carga Nivel)};

    % Módulo LoRa SPI
    \draw[thick, fill=orange!5] (0,6.0) rectangle (6.5,7.5);
    \node at (3,6.9) {\textbf{LoRa Ra-02 (SPI1)}};
    \node[right] at (0.2,6.3) {\footnotesize SCK:10 | MOSI:11 | MISO:12 | CS:13};

    % Interconexión RP2040 - LoRa
    \draw[<->, thick] (2.25,5.2) -- (2.25,6.0);

    % Bloque de Aislamiento y Potencia
    \draw[thick, fill=green!5] (6.5,0.2) rectangle (11.5,4.5);
    \node at (9,4.0) {\textbf{Etapa de Aislamiento}};
    \node at (9,3.4) {\textbf{\& Drivers de Potencia}};
    \node[right] at (7,2.5) {\footnotesize IN1 (Relé TDS)};
    \node[right] at (7,1.0) {\footnotesize IN2 (Relé Nivel)};

    % Conexiones Lógicas RP2040 -> Potencia
    \draw[->, thick, color=blue!80!black] (4.5,1.4) -- (6.5,2.2);
    \draw[->, thick, color=blue!80!black] (4.5,0.6) -- (6.5,1.0);

    % Cargas en 110V AC
    \draw[thick, fill=red!5] (12.0,2.0) rectangle (15.2,3.2);
    \node at (13.6,2.6) {\small Electroválvula};
    \node[below] at (13.6,2.1) {\footnotesize (110V AC)};

    \draw[thick, fill=red!5] (12.0,0.4) rectangle (15.2,1.6);
    \node at (13.6,1.0) {\small Bomba Recirc.};
    \node[below] at (13.6,0.5) {\footnotesize (110V AC)};

    % Salidas hacia cargas AC
    \draw[->, thick, color=red!70!black] (10.0,2.6) -- (12.0,2.6);
    \draw[->, thick, color=red!70!black] (10.0,1.0) -- (12.0,1.0);
\end{circuitikz}
\caption{Diagrama de bloques simplificado para el nodo RP2040, interfaz LoRa y etapa de potencia aislada.}
\label{fig:esquematico}
\end{figure}

\begin{table}[H]
\centering
\caption{Asignación de Pines del Microcontrolador RP2040}
\vspace{2mm}
\begin{tabular}{llll}
\toprule
\textbf{Pin RP2040} & \textbf{Periférico} & \textbf{Dispositivo / Señal} & \textbf{Tipo de Interfaz} \\ \midrule
GPIO26 (ADC0) & ADC & Sensor de SDT / Conductividad & Analógica (0--3.3V) \\
GPIO15 & Output & Trigger Sensor Ultrasónico JSN-SR04T & Digital (Salida) \\
GPIO14 & Input & Echo Sensor Ultrasónico JSN-SR04T & Digital (Entrada) \\
GPIO16 & Output & Control Optoacoplador Carga TDS & Digital (Salida) \\
GPIO17 & Output & Control Optoacoplador Carga Nivel & Digital (Salida) \\
GPIO10--GPIO13 & SPI1 & LoRa Ra-02 (SCK, MOSI, MISO, CS) & Bus SPI \\
GPIO27, GPIO28 & GPIO & LoRa DIO0 e Interrupción / Reset & Control Digital \\ \bottomrule
\end{tabular}
\end{table}

\section{Protocolo de Comunicación y Estructura MQTT}
La comunicación entre el nodo RP2040 y el Gateway ESP32 utiliza tramas LoRa cifradas con AES-128 (modo ECB). Posteriormente, el Gateway publica la información procesada hacia el broker Mosquitto en la Raspberry Pi 5.

\begin{table}[H]
\centering
\caption{Estructura de Tópicos MQTT y Payloads}
\vspace{2mm}
\begin{tabular}{lp{5.5cm}p{5.5cm}}
\toprule
\textbf{Tópico MQTT} & \textbf{Formato del Payload} & \textbf{Descripción y Uso} \\ \midrule
\texttt{sensor/data/out} & \texttt{\{"read\_time": "...", "avg\_tds\_ppm": 4.5, "avg\_depth\_cm": 280.2, "tds\_load": 1, "depth\_load": 0\}} & Publicación del promedio acumulado (bloque de 20 muestras) o lecturas instantáneas. \\ \midrule
\texttt{sensor/data/in} & \texttt{"tds/load/on"}, \texttt{"tds/load/off"}, \texttt{"depth/load/on"}, o JSON & Suscripción a comandos remotos recibidos desde la interfaz Web/Grafana. \\ \bottomrule
\end{tabular}
\end{table}

\section{Implementación de Software}
A continuación se presentan los módulos principales de software implementados en MicroPython para el nodo perimetral RP2040 y el Gateway ESP32.

\subsection{Firmware del Nodo Perimetral (RP2040)}
El código cargado en la Raspberry Pi Pico realiza mediciones no bloqueantes, aplica cifrado AES-128 y conmuta las cargas en respuesta a comandos recibidos por LoRa.

\begin{lstlisting}[language=Python,caption={Firmware completo del nodo sensor y actuador en RP2040 (\texttt{main.py})},label={lst:rp2040_py}]
from machine import Pin, SPI, time_pulse_us, ADC, idle
from lora_lib import LoRa
import time
import ujson
import ucryptolib
import ubinascii

# ==========================================
# CONFIGURACIÓN CRIPTOGRÁFICA (AES-128 ECB)
# ==========================================
AES_KEY = b"MiClaveSecreta16"

def pad(texto):
    pad_len = 16 - (len(texto) % 16)
    return texto + bytes([pad_len]) * pad_len

def unpad(texto):
    pad_len = texto[-1]
    if pad_len < 1 or pad_len > 16:
        raise ValueError("Padding invalido")
    return texto[:-pad_len]

def cifrar_mensaje(mensaje_str):
    cipher = ucryptolib.aes(AES_KEY, 1)  # 1 = ECB
    cifrado_bytes = cipher.encrypt(pad(mensaje_str.encode('utf-8')))
    return ubinascii.hexlify(cifrado_bytes).decode('utf-8')

def descifrar_mensaje(mensaje_hex_str):
    try:
        cifrado_bytes = ubinascii.unhexlify(mensaje_hex_str)
        cipher = ucryptolib.aes(AES_KEY, 1)
        return unpad(cipher.decrypt(cifrado_bytes)).decode('utf-8')
    except Exception as e:
        print(f"Error descifrando: {e}")
        return None

# ==========================================
# CONFIGURACIÓN SENSORES Y CARGAS
# ==========================================
trig = Pin(15, Pin.OUT)
echo = Pin(14, Pin.IN)
tds_sensor = ADC(Pin(26))

carga_tds = Pin(16, Pin.OUT)
carga_nivel = Pin(17, Pin.OUT)

def medir_distancia(temperatura_c=20.0):
    velocidad_cm_us = (331.3 + (0.6 * temperatura_c)) * 0.0001
    trig.low(); time.sleep_us(5)
    trig.high(); time.sleep_us(15); trig.low()
    duracion = time_pulse_us(echo, 1, 30000)
    if duracion > 0:
        return (duracion * velocidad_cm_us) / 2.0
    return -1

def calcular_tds(valor_adc):
    voltage = (valor_adc / 65535.0) * 3.3
    tds_value = (133.42 * (voltage ** 3) - 255.86 * (voltage ** 2) + 857.39 * voltage) * 0.5
    return max(0.0, tds_value)

def timestamp_ms():
    t = time.localtime()
    try:
        ms = (time.time_ns() // 1000000) % 1000
    except Exception:
        ms = time.ticks_ms() % 1000
    return f"{t[0]:04d}-{t[1]:02d}-{t[2]:02d} {t[3]:02d}:{t[4]:02d}:{t[5]:02d}.{ms:03d}"

# ==========================================
# ESTADO NO BLOQUEANTE
# ==========================================
SEND_INTERVAL_MS = 3000
TX_GUARD_MS = 200

rx_queue = []
rx_resume_at = None
ultimo_envio_datos = time.ticks_add(time.ticks_ms(), -SEND_INTERVAL_MS)
datos_exitosos = 0

def al_recibir(dato):
    if len(rx_queue) < 8:
        rx_queue.append(dato)

def transmitir(payload_hex):
    global rx_resume_at
    lora.send(payload_hex)
    rx_resume_at = time.ticks_add(time.ticks_ms(), TX_GUARD_MS)

# ==========================================
# INICIALIZACIÓN LORA
# ==========================================
spi = SPI(1, baudrate=1000000, sck=Pin(10), mosi=Pin(11), miso=Pin(12))
lora = LoRa(spi, cs=13, reset=28, dio0=27, callback=al_recibir)
lora.set_bandwidth(250000)
lora.set_rx_mode()

print("--- INICIANDO SISTEMA RP2040 CIFRADO ---")

def procesar_rx():
    global datos_exitosos
    while rx_queue:
        dato_crudo = rx_queue.pop(0)
        try:
            texto_hex = dato_crudo.decode('utf-8') if isinstance(dato_crudo, (bytes, bytearray)) else str(dato_crudo)
            texto_descifrado = descifrar_mensaje(texto_hex)
            if not texto_descifrado:
                continue

            msg = ujson.loads(texto_descifrado)
            tipo = msg.get("tipo")
            if tipo == "ack":
                datos_exitosos += 1
            elif tipo == "comando":
                valor_tds = msg.get("tds_load", 0)
                valor_nivel = msg.get("depth_load", 0)
                carga_tds.value(valor_tds)
                carga_nivel.value(valor_nivel)
        except Exception as e:
            print("Error procesando trama entrante:", e)

def enviar_datos():
    dist = medir_distancia()
    tds_ppm = calcular_tds(tds_sensor.read_u16())
    datos = f"{timestamp_ms()} {round(tds_ppm, 1)} {round(dist, 1) if dist >= 0 else None}"
    payload_cifrado = cifrar_mensaje(datos)
    transmitir(payload_cifrado)

while True:
    ahora = time.ticks_ms()
    if rx_resume_at is not None and time.ticks_diff(ahora, rx_resume_at) >= 0:
        lora.set_rx_mode()
        rx_resume_at = None

    procesar_rx()

    if rx_resume_at is None and time.ticks_diff(ahora, ultimo_envio_datos) >= SEND_INTERVAL_MS:
        ultimo_envio_datos = ahora
        enviar_datos()

    idle()
\end{lstlisting}

\subsection{Firmware del Gateway (ESP32)}
El Gateway puentea los paquetes LoRa cifrados hacia la red Wi-Fi mediante MQTT cifrado sobre TLS, acumulando bloques de 20 muestras antes de emitir decisiones de control local.

\begin{lstlisting}[language=Python,caption={Firmware del Gateway ESP32 (Puente LoRa--MQTT)},label={lst:esp32_py}]
import network, time, json, ucryptolib, ubinascii
from machine import Pin, SPI
from umqtt.simple import MQTTClient
from lora_lib import LoRa

WIFI_SSID = "IoT"
WIFI_PASS = "EIJabcd1234"
MQTT_BROKER = "192.168.0.115"
MQTT_PORT = 8883
MQTT_CLIENT_ID = "esp32smv"
PUB_TOPIC = b"sensor/data/out"
SUB_TOPIC = b"sensor/data/in"

AES_KEY = b"MiClaveSecreta16"

def unpad(data):
    return data[:-data[-1]]

def encrypt_message(msg_str):
    cipher = ucryptolib.aes(AES_KEY, 1)
    pad_len = 16 - (len(msg_str) % 16)
    padded = msg_str.encode('utf-8') + bytes([pad_len]) * pad_len
    return ubinascii.hexlify(cipher.encrypt(padded)).decode('utf-8')

def decrypt_message(hex_str):
    try:
        cipher = ucryptolib.aes(AES_KEY, 1)
        raw = ubinascii.unhexlify(hex_str.strip())
        return unpad(cipher.decrypt(raw)).decode('utf-8')
    except Exception as e:
        return None

MAX_MUESTRAS = 20
historial_tds, historial_depth = [], []
cmd_tds, cmd_depth = 0, 0

def enviar_comando_lora():
    global cmd_tds, cmd_depth
    payload = json.dumps({"tipo": "comando", "tds_load": cmd_tds, "depth_load": cmd_depth})
    time.sleep_ms(100)
    lora.send(encrypt_message(payload))
    lora.set_rx_mode()

def connect_wifi():
    wlan = network.WLAN(network.STA_IF)
    wlan.active(True)
    if not wlan.isconnected():
        wlan.connect(WIFI_SSID, WIFI_PASS)
        while not wlan.isconnected():
            time.sleep(0.5)

def lora_receive(msg):
    global cmd_tds, cmd_depth, historial_tds, historial_depth
    plain = decrypt_message(msg)
    if plain is None: return
    
    date, _time, tds_str, depth_str = plain.split()
    historial_tds.append(float(tds_str))
    historial_depth.append(float(depth_str))
    
    if len(historial_tds) >= MAX_MUESTRAS:
        prom_tds = sum(historial_tds) / MAX_MUESTRAS
        prom_depth = sum(historial_depth) / MAX_MUESTRAS
        cmd_tds = 1 if prom_tds <= 5.0 else 0
        cmd_depth = 1 if prom_depth <= 300.0 else 0
        
        data_mqtt = {"read_time": f"{date} {_time}", "avg_tds_ppm": round(prom_tds,2), "avg_depth_cm": round(prom_depth,2), "tds_load": cmd_tds, "depth_load": cmd_depth}
        client.publish(PUB_TOPIC, json.dumps(data_mqtt))
        enviar_comando_lora()
        historial_tds.clear(); historial_depth.clear()
    else:
        enviar_comando_lora()

connect_wifi()
client = MQTTClient(MQTT_CLIENT_ID, MQTT_BROKER, port=MQTT_PORT, ssl=True, ssl_params={"server_hostname": "localhost"})
client.connect()
client.subscribe(SUB_TOPIC)

spi = SPI(2, baudrate=1000000, sck=Pin(14), mosi=Pin(13), miso=Pin(12))
lora = LoRa(spi, cs=17, reset=5, dio0=6, callback=lora_receive)
lora.set_rx_mode()

while True:
    client.check_msg()
    time.sleep(0.1)
\end{lstlisting}

\section{Conclusiones y Recomendaciones}
\begin{enumerate}[1.]
    \item La migración hacia el procesador RP2040 en los nodos perimetrales proporciona un manejo determinista y preciso para los sensores de campo, mientras que la Raspberry Pi 5 ofrece la capacidad de cómputo necesaria para gestionar la pila de almacenamiento y análisis de datos sin cuellos de botella.
    \item Es imperativo mantener el aislamiento galvánico mediante contactores/relés de alta calidad para evitar interferencias electromagnéticas (EMI) durante el arranque del motor de la bomba de recirculación de 110V AC.
\end{enumerate}

\end{document}